\section{Detection}\label{sec:detection}

The evolution of object detection paradigms in our framework moves from classical computer vision to domain-adapted deep learning.
As an initial baseline, a Mixture of Gaussians (MOG2) background subtraction approach \cite{survey} enhanced with CLAHE preprocessing \cite{clahe} proves highly vulnerable to court floor specularities, player stasis, and dynamic audience backgrounds.
Transitioning to an object detection paradigm via a COCO-pre-trained YOLOv11 model \cite{yolo_paper} substantially mitigates these spatial errors by extracting robust semantic features.
However, significant domain gaps persist with the off-the-shelf model: the basketball is entirely absent from the pre-trained taxonomy, and players and referees are uniformly grouped under a single ``person'' class, precluding downstream identity distinction and necessitating a fine-tuned, domain-specific architecture.

% This section describes the three detection approaches evaluated in this work, progressing from a classical background subtraction baseline to a fine-tuned deep learning model with multi-pass inference.

% \subsection{Background Subtraction (MOG2)}
% As a baseline, we employ Mixture of Gaussians (MOG2) background subtraction \cite{survey}.
% To mitigate illumination variations, frames are preprocessed using Contrast Limited Adaptive Histogram Equalization (CLAHE) \cite{clahe}.
% The generated foreground mask is subsequently refined via morphological operations and connected component analysis.

% Despite these enhancements, this baseline suffers from four primary failure modes.
% First, floor specularities create intense reflections and shadows that evade MOG2's color-based shadow detection.
% Then, temporarily stationary players are erroneously integrated into the background model, and dynamic backgrounds, such as moving spectators, introduce persistent false positives.

% \subsection{Pre-Trained YOLOv11}
% To overcome MOG2 limitations, we deploy a YOLOv11 architecture \cite{yolo_paper} pre-trained on the COCO dataset.
% Transitioning to this object detection paradigm substantially improves accuracy and mitigates the limitations of the background subtraction baseline.

% However, two distinct domain gaps persist: the sports ball is entirely absent from the pre-trained taxonomy, and players and referees are uniformly classified as ``person'', precluding downstream identity distinction.

\subsection{Fine-Tuned YOLOv26 with Multi-Pass Inference}
To bridge these domain gaps, we adopt the YOLOv26 architecture and fine-tune its backbone on a domain-specific annotated dataset.
The model is trained to recognise each player individually, as well as the ball, so identity is resolved directly by the detector.

However, because the ball occupies a minimal pixel area, standard single-pass inference downsamples it too aggressively for reliable localization.
We circumvent this spatial resolution trade-off via a dual-pass inference strategy: a primary pass optimized for player detection, and a parallel high-resolution pass configured with a permissive confidence threshold specifically tailored to isolate the ball.
The resulting detections are subsequently unified into a single composite output.

Finally, to eliminate duplicate bounding boxes when a player is detected under multiple conflicting identities, we substitute standard non-maximum suppression with a class-agnostic variant \cite{nms}.
This suppresses redundant detections by spatial overlap alone, ensuring a single clean bounding box per player is passed to the downstream tracker.
 
% \subsubsection{Class-Agnostic NMS}
% A player may be detected multiple times under different identities, which ordinary non-maximum suppression cannot merge since it only compares same-identity boxes.
% We instead apply a class-agnostic variant \cite{nms}, suppressing by spatial overlap alone to keep one box per player, yielding the final detections passed to the tracker.

\begin{figure}[h]
  \centering
  \includegraphics[width=\linewidth]{../media/fig_detection.png}
  \caption{Sample detection output from fine-tuned YOLOv26 with class-agnostic NMS
           on cam\_13 (three selected frames).}
  \label{fig:detection}
\end{figure}
